\documentclass[fontset=windows, chinese, 10pt, xcolor=dvipsnames]{beamer}

\usepackage{amsmath,amssymb,mathtools}
\usepackage{xeCJK}
\usepackage{zhnumber}
\usepackage{tikz}
\usetikzlibrary{quotes,angles,arrows.meta,calc}

\usetheme{metropolis}
\usepackage{beamercolorthemeHistNanjing}

\newcommand{\red}[1]{\textcolor{red}{#1}}
\newcommand{\blue}[1]{\textcolor{blue}{#1}}
\newcommand{\green}[1]{\textcolor{green}{#1}}
\newcommand{\gray}[1]{\textcolor{gray}{#1}}
\newcommand{\orange}[1]{\textcolor{orange}{#1}}

\newcommand{\vecfig}[1]{%
  \IfFileExists{tikz/#1.tex}%
    {\input{tikz/#1.tex}}%
    {\input{../../lecture_new/chapters/vec/tikz/#1.tex}}}

\newcommand{\sssub}[1]{ {\mathrm{#1} }}  % script size subscript\\\
\newcommand{\funcd}[2]{ \frac{d{#1}}{d{#2}}} %function full derivative
\newcommand{\funcpd}[2]{ \frac{\partial{#1}}{\partial{#2}}} %function partial derivative
\newcommand{\fnald}[2] { \frac{\delta {#1}}{\delta {#2}}} % functional derivative

\renewcommand{\Re}{\operatorname{Re}}
\renewcommand{\Im}{\operatorname{Im}}
\def\ext{\sssub{ext}}
\def\x{\sssub{x}}
\def\xc{\sssub{xc}}
\def\H{\sssub{H}}
\def\Hx{\sssub{Hx}}
\def\Hxc{\sssub{Hxc}}
\def\c{\sssub{c}}
\def\s{\sssub{s}}

% \def\ext{\sssub{ext}}
% \def\ext{\sssub{ext}}
% \def\ext{\sssub{ext}}
% \def\ext{\sssub{ext}}
% \defsssub{ext}


% definitions of bold faced symbols
\def\brpppp{{\mathbf{r}^{\prime\prime\prime\prime}}}
\def\brppp{{\mathbf{r}^{\prime\prime\prime}}}
\def\brpp{{\mathbf{r}^{\prime\prime}}}
\def\brp{{\mathbf{r}^{\prime}}}
\def\bzp{{\mathbf{z}^{\prime}}}
\def\bxp{{\mathbf{x}^{\prime}}}
\def\tp{{{t}^{\prime}}}
\def\tpp{{{t}^{\prime\prime}}}
\def\tppp{{{t}^{\prime\prime\prime}}}

\def\tbr{{\tilde{\mathbf{r}}}}
\def\bk{{\mathbf{k}}}
\def\br{{\mathbf{r}}}
\def\bp{{\mathbf{p}}}
\def\bq{{\mathbf{q}}}
\def\bv{{\mathbf{v}}}
\def\bz{{\mathbf{z}}}
\def\mbf{{\mathbf{f}}}

\def\bx{{\mathbf{x}}}
\def\bR{{\mathbf{R}}}
\def\bM{{\mathbf{M}}}
\def\bP{{\mathbf{P}}}
\def\bT{{\mathbf{T}}}
\def\bK{{\mathbf{K}}}
\def\bA{{\mathbf{A}}}
\def\bB{{\mathbf{B}}}
\def\bD{{\mathbf{D}}}
\def\bE{{\mathbf{E}}}
\def\bF{{\mathbf{F}}}
\def\bG{{\mathbf{G}}}
\def\bH{{\mathbf{H}}}
\def\bI{{\mathbf{I}}}
\def\bJ{{\mathbf{J}}}
\def\bP{{\mathbf{P}}}
\def\bS{{\mathbf{S}}}
\def\bV{{\mathbf{V}}}
\def\bX{{\mathbf{X}}}
\def\bY{{\mathbf{Y}}}
\def\dd{{\mathop{}\!\mathrm{d}}}
\def\dr{{\!\!\dd\br\,}}
\def\dx{{\!\!\dd\bx\,}}
\def\dy{{\!\!\dd\by\,}}
\def\dk{{\!\!\dd\bk\,}}
\def\dG{{\!\!\dd\bG\,}}


\def\half{{\frac{1}{2}}}
\def\tr{{\mathrm{tr}}}


\def\Arg{{\mathrm{Arg}\;}}
\def\arg{{\mathrm{arg}\;}}
\def\Res{{\mathrm{Res}\;}}


\title{矢量分析}
\subtitle{第一讲：矢量运算}


\date{\zhdigits{\the\year},\zhganzhinian{\the\year}年秋}
\author{罗凯}
\institute{南京理工大学物理学院}

\usebackgroundtemplate{%
\begin{tikzpicture}[remember picture,overlay]
  \node[opacity=0.05, at=(current page.center)] {
    \includegraphics[width=0.6\paperwidth]{njust-logo.jpeg}
  };
\end{tikzpicture}}
\begin{document}

\begin{frame}
  \titlepage
\end{frame}

\AtBeginSubsection{
  \begin{frame}{提纲}
    \tableofcontents[currentsection,currentsubsection]
  \end{frame}
}

\begin{frame}
  \titlepage
\end{frame}

\AtBeginSubsection{
  \begin{frame}{提纲}
    \tableofcontents[currentsection,currentsubsection]
  \end{frame}
}

\section{矢量}

\begin{frame}{本章要抓住的那根纲}
  \begin{quote}
    引其纲，万目皆张。\\
    \hfill ——《吕氏春秋·用民》
  \end{quote}
  渔网的总绳一旦被拉住，所有网眼就跟着张开。
  矢量就是后面矩阵、微分算子和张量的这根总绳。
  \vspace{0.4em}
  同一套材料的标准叙述见 Arfken、Weber 与 Harris，以及 Byron 与 Fuller。
\end{frame}

\begin{frame}{标量与矢量}
  \textbf{矢量}既有大小，又有方向。
  \textbf{标量}只有大小。
  \vspace{0.5em}
  物理里常见的矢量：
  \begin{itemize}
    \item 力学中的力、速度、加速度；
    \item 电磁学中的电磁场；
    \item 流体力学中的流动速度。
  \end{itemize}
  维数可以高于三维，从而进入更一般的高维空间。
  这些空间之间的线性映射用矩阵表示：当 $A\mathbf{x}=\mathbf{b}$ 可逆时，
  \begin{equation*}
    \mathbf{x}=A^{-1}\mathbf{b}.
  \end{equation*}
  逆矩阵的构造放在第二讲。
  \vspace{0.3em}
  本章把矢量写成粗体，例如 $\mathbf{A}$。
\end{frame}

\subsection{分量与运算}

\begin{frame}{直角分量}
  \begin{columns}[c]
    \begin{column}{0.52\textwidth}
      三维直角坐标系中
      \begin{equation*}
        \mathbf{A}
        = A_x\mathbf{e}_x
        + A_y\mathbf{e}_y
        + A_z\mathbf{e}_z.
      \end{equation*}
      $\{\mathbf{e}_x,\mathbf{e}_y,\mathbf{e}_z\}$ 是单位正交基。
      两个矢量相等，当且仅当每一对对应分量都相等。
    \end{column}
    \begin{column}{0.46\textwidth}
      \centering
      \vecfig{components}
    \end{column}
  \end{columns}
\end{frame}

\begin{frame}{模}
  由勾股定理，
  \begin{equation*}
    |\mathbf{A}|
    = \sqrt{A_x^2+A_y^2+A_z^2}.
  \end{equation*}
  只在 $xy$ 平面内时 $A_z=0$，公式退化为
  $\sqrt{A_x^2+A_y^2}$。
  \vspace{0.6em}
  零矢量的每个分量都是 $0$，模也是 $0$。
  它是唯一没有方向的矢量。
\end{frame}

\begin{frame}{标量与数乘}
  温度、质量、电荷都由一个大小描写，所以是标量。
  用标量 $\alpha$ 去乘矢量，就是逐个分量去乘：
  \begin{equation*}
    \alpha\mathbf{A}
    = \alpha A_x\mathbf{e}_x
    + \alpha A_y\mathbf{e}_y
    + \alpha A_z\mathbf{e}_z.
  \end{equation*}
  因此
  \begin{equation*}
    |\alpha\mathbf{A}|
    = \sqrt{\alpha^2(A_x^2+A_y^2+A_z^2)}
    = |\alpha|\,|\mathbf{A}|.
  \end{equation*}
  \begin{itemize}
    \item $\alpha>0$：方向不变，长度乘以 $\alpha$。
    \item $\alpha<0$：方向反转，长度乘以 $|\alpha|$。
    \item $\alpha=0$：得到零矢量。
  \end{itemize}
\end{frame}

\begin{frame}{加法与减法}
  加减法逐分量进行：
  \begin{align*}
    \mathbf{A}\pm\mathbf{B}
    &= (A_x\pm B_x)\mathbf{e}_x
     + (A_y\pm B_y)\mathbf{e}_y \\
    &\quad + (A_z\pm B_z)\mathbf{e}_z.
  \end{align*}
  几何上，加法服从平行四边形法则。
  把两矢量首尾相接时，$\mathbf{B}-\mathbf{A}$ 是从 $\mathbf{A}$ 的箭尖指向 $\mathbf{B}$ 的箭尖的矢量。
\end{frame}

\begin{frame}{加法与减法的图}
  \centering
  \scalebox{0.92}{\vecfig{sum_diff}}
  \vspace{0.4em}

  {\small 图中 $\mathbf{A}=4\mathbf{e}_x+4\mathbf{e}_y$，
  $\mathbf{B}=\mathbf{e}_x+8\mathbf{e}_y$。
  下一页把这个例子算完。}
\end{frame}

\begin{frame}{例：和与差}
  取 $\mathbf{A}=4\mathbf{e}_x+4\mathbf{e}_y$，
  $\mathbf{B}=\mathbf{e}_x+8\mathbf{e}_y$。
  逐分量相加：
  \begin{align*}
    \mathbf{S}
    &= \mathbf{A}+\mathbf{B}
     = (4+1)\mathbf{e}_x+(4+8)\mathbf{e}_y \\
    &= 5\mathbf{e}_x+12\mathbf{e}_y.
  \end{align*}
  模为
  \begin{align*}
    |\mathbf{S}|
    &= \sqrt{5^2+12^2}
     = \sqrt{25+144}
     = \sqrt{169}
     = 13.
  \end{align*}
\end{frame}

\begin{frame}{例：和与差（续）}
  差为
  \begin{align*}
    \mathbf{D}
    &= \mathbf{B}-\mathbf{A}
     = (1-4)\mathbf{e}_x+(8-4)\mathbf{e}_y \\
    &= -3\mathbf{e}_x+4\mathbf{e}_y.
  \end{align*}
  因而
  \begin{align*}
    |\mathbf{D}|
    &= \sqrt{(-3)^2+4^2}
     = \sqrt{9+16}
     = \sqrt{25}
     = 5.
  \end{align*}
  $5$--$12$--$13$ 与 $3$--$4$--$5$ 两组勾股数正好核对这两个模。
\end{frame}

\subsection{点积}

\begin{frame}{点积的两种写法}
  两个矢量的点积（标量积）是一个标量：
  \begin{equation*}
    \mathbf{A}\cdot\mathbf{B}
    = |\mathbf{A}|\,|\mathbf{B}|\cos\theta,
  \end{equation*}
  其中 $\theta$ 是二者的夹角。
  在直角分量下，利用
  $\mathbf{e}_i\cdot\mathbf{e}_j=\delta_{ij}$ 展开：
  \begin{align*}
    \mathbf{A}\cdot\mathbf{B}
    &= (A_x\mathbf{e}_x+A_y\mathbf{e}_y+A_z\mathbf{e}_z)
      \cdot
      (B_x\mathbf{e}_x+B_y\mathbf{e}_y+B_z\mathbf{e}_z) \\
    &= A_x B_x+A_y B_y+A_z B_z.
  \end{align*}
  交叉项因为基矢正交而消失。
\end{frame}

\begin{frame}{点积读出夹角}
  \begin{columns}[c]
    \begin{column}{0.48\textwidth}
      由 $\cos\theta$ 的符号：
      \begin{itemize}
        \item $\mathbf{A}\cdot\mathbf{B}>0$：锐角；
        \item $\mathbf{A}\cdot\mathbf{B}=0$：两矢量垂直；
        \item $\mathbf{A}\cdot\mathbf{B}<0$：钝角。
      \end{itemize}
      特别地，$\mathbf{A}\cdot\mathbf{A}=|\mathbf{A}|^2$。
    \end{column}
    \begin{column}{0.50\textwidth}
      \centering
      \scalebox{0.82}{\vecfig{dot_product}}
    \end{column}
  \end{columns}
\end{frame}

\begin{frame}{例：由点积求夹角}
  设
  $\mathbf{A}=\sqrt{3}\,\mathbf{e}_x+\mathbf{e}_y$，
  $\mathbf{B}=2\mathbf{e}_x$。
  先算模：
  \begin{align*}
    |\mathbf{A}|
    &= \sqrt{(\sqrt{3})^2+1^2}
     = \sqrt{3+1}
     = \sqrt{4}
     = 2,\\
    |\mathbf{B}|
    &= \sqrt{2^2}=2.
  \end{align*}
  再算点积：
  \begin{equation*}
    \mathbf{A}\cdot\mathbf{B}
    = (\sqrt{3})(2)+1\cdot 0
    = 2\sqrt{3}.
  \end{equation*}
\end{frame}

\begin{frame}{例：由点积求夹角（续）}
  代入几何定义
  $\mathbf{A}\cdot\mathbf{B}=|\mathbf{A}|\,|\mathbf{B}|\cos\theta$：
  \begin{equation*}
    2\sqrt{3}
    = 2\cdot 2\cos\theta
    = 4\cos\theta,
  \end{equation*}
  所以
  \begin{equation*}
    \cos\theta=\frac{\sqrt{3}}{2},
    \qquad
    \theta
    = \cos^{-1}\!\Bigl(\frac{\sqrt{3}}{2}\Bigr)
    = 30^\circ.
  \end{equation*}
  $\mathbf{B}$ 沿 $x$ 轴，$\mathbf{A}$ 与它成 $30^\circ$，与前图一致。
\end{frame}

\begin{frame}{例：水分子的键角}
  % \centering
  \begin{figure}
  \centering
  % \includegraphics[width=0.8\textwidth]{water_bond.png}
  \vecfig{water_bond}
  \end{figure}

  把氧原子放在原点，两个氢原子（单位：\AA）为
  \begin{equation*}
    \mathrm{H}_1=(0.757,\,0,\,0.586),
    \qquad
    \mathrm{H}_2=(-0.757,\,0,\,0.586).
  \end{equation*}
  键矢量 $\mathbf{u}=\overrightarrow{\mathrm{OH}_1}$，
  $\mathbf{v}=\overrightarrow{\mathrm{OH}_2}$。
  记 $a=0.757$，$b=0.586$，则
  \begin{equation*}
    \mathbf{u}=(a,0,b),
    \qquad
    \mathbf{v}=(-a,0,b).
  \end{equation*}
\end{frame}

\begin{frame}{例：水分子的键角}
  \begin{align*}
    \mathbf{u}\cdot\mathbf{v}
    &= a(-a)+0\cdot 0+b\cdot b
     = b^2-a^2,\\
    |\mathbf{u}|^2
    &= a^2+b^2
     = |\mathbf{v}|^2.
  \end{align*}
  因此
  \begin{equation*}
    \cos\alpha
    = \frac{\mathbf{u}\cdot\mathbf{v}}{|\mathbf{u}|\,|\mathbf{v}|}
    = \frac{b^2-a^2}{a^2+b^2}.
  \end{equation*}
  代入数值：
  \begin{align*}
    a^2&=0.757^2=0.573049,\\
    b^2&=0.586^2=0.343396,\\
    \cos\alpha
    &= \frac{0.343396-0.573049}{0.573049+0.343396}
     = \frac{-0.229653}{0.916445}
     \approx -0.251.
  \end{align*}
  所以 $\alpha=\arccos(-0.251)\approx 104.5^\circ$。
\end{frame}

% \begin{frame}{例：水分子的键角（图）}
%   \centering
%   \scalebox{0.62}{\vecfig{water_bond}}
% \end{frame}

\begin{frame}{例：用点积判断垂直}
  判断
  $\mathbf{p}=2\mathbf{e}_x-\mathbf{e}_y+2\mathbf{e}_z$
  与
  $\mathbf{q}=\mathbf{e}_x+2\mathbf{e}_y$
  是否垂直。
  \vspace{0.5em}
  \begin{align*}
    \mathbf{p}\cdot\mathbf{q}
    &= 2\cdot 1+(-1)\cdot 2+2\cdot 0 \\
    &= 2-2+0
     = 0.
  \end{align*}
  点积为零，所以 $\mathbf{p}\perp\mathbf{q}$。
  一旦点积已经是零，就不必再算两个模。
\end{frame}

\begin{frame}{小结}
  \begin{itemize}
    \item $\mathbf{A}=A_x\mathbf{e}_x+A_y\mathbf{e}_y+A_z\mathbf{e}_z$，
      $|\mathbf{A}|=\sqrt{A_x^2+A_y^2+A_z^2}$。
    \item 数乘逐分量进行；负标量使方向反转。
    \item 加法是平行四边形法则；$\mathbf{B}-\mathbf{A}$ 从 $\mathbf{A}$ 的箭尖指向 $\mathbf{B}$ 的箭尖。
    \item $\mathbf{A}\cdot\mathbf{B}=|\mathbf{A}|\,|\mathbf{B}|\cos\theta
      =A_xB_x+A_yB_y+A_zB_z$，
      后一式来自基矢的正交性。
    \item 点积为零就是垂直的判据。
  \end{itemize}
\end{frame}

\section{矢量的乘积}

\subsection{叉积}

\begin{frame}{叉积}
  $\mathbf{A}\times\mathbf{B}$ 同时垂直于 $\mathbf{A}$ 和 $\mathbf{B}$。
  它的模等于二者所张平行四边形的面积：
  \begin{equation*}
    |\mathbf{A}\times\mathbf{B}|
    = |\mathbf{A}|\,|\mathbf{B}|\sin\theta.
  \end{equation*}
  方向由右手定则确定。
  分量形式写成行列式
  \begin{equation*}
    \mathbf{A}\times\mathbf{B}
    =
    \begin{vmatrix}
      \mathbf{e}_x & \mathbf{e}_y & \mathbf{e}_z \\
      A_x & A_y & A_z \\
      B_x & B_y & B_z
    \end{vmatrix}.
  \end{equation*}
\end{frame}

\begin{frame}{按第一行展开}
  把行列式按第一行展开：
  \begin{align*}
    \mathbf{A}\times\mathbf{B}
    &= \mathbf{e}_x(A_y B_z-A_z B_y)
     - \mathbf{e}_y(A_x B_z-A_z B_x) \\
    &\quad + \mathbf{e}_z(A_x B_y-A_y B_x).
  \end{align*}
  由此直接读出三条性质：
  \begin{itemize}
    \item 交换 $\mathbf{A}$ 与 $\mathbf{B}$ 相当于交换两行，行列式变号，所以
      $\mathbf{A}\times\mathbf{B}=-\mathbf{B}\times\mathbf{A}$。
    \item $\mathbf{A}=\mathbf{B}$ 时两行相同，所以 $\mathbf{A}\times\mathbf{A}=\mathbf{0}$。
    \item 两矢量平行时 $\sin\theta=0$，叉积为零。
  \end{itemize}
\end{frame}

\begin{frame}{叉积的几何}
  \begin{columns}[c]
    \begin{column}{0.40\textwidth}
      叉积立在平行四边形所在平面之外，长度等于这块面积，指向由右手定则给出。
    \end{column}
    \begin{column}{0.58\textwidth}
      \centering
      \scalebox{0.9}{\vecfig{cross_product}}
    \end{column}
  \end{columns}
\end{frame}

\begin{frame}{例：把行列式展开}
  设
  $\mathbf{A}=-\mathbf{e}_x+\mathbf{e}_y+\mathbf{e}_z$，
  $\mathbf{B}=\mathbf{e}_x-\mathbf{e}_y+\mathbf{e}_z$。
  \begin{align*}
    \mathbf{A}\times\mathbf{B}
    &=
    \begin{vmatrix}
      \mathbf{e}_x & \mathbf{e}_y & \mathbf{e}_z \\
      -1 & 1 & 1 \\
      1 & -1 & 1
    \end{vmatrix} \\
    &= \mathbf{e}_x\bigl(1\cdot 1-1\cdot(-1)\bigr)
     - \mathbf{e}_y\bigl((-1)\cdot 1-1\cdot 1\bigr) \\
    &\quad + \mathbf{e}_z\bigl((-1)\cdot(-1)-1\cdot 1\bigr) \\
    &= \mathbf{e}_x(1+1)
     - \mathbf{e}_y(-1-1)
     + \mathbf{e}_z(1-1) \\
    &= 2\mathbf{e}_x+2\mathbf{e}_y.
  \end{align*}
\end{frame}

\begin{frame}{例：单位法向}
  上一步的模是
  \begin{equation*}
    |\mathbf{A}\times\mathbf{B}|
    = \sqrt{2^2+2^2}
    = \sqrt{8}
    = 2\sqrt{2}.
  \end{equation*}
  右手定向的单位法向为
  \begin{equation*}
    \mathbf{e}_n
    = \frac{\mathbf{A}\times\mathbf{B}}{|\mathbf{A}\times\mathbf{B}|}
    = \frac{2\mathbf{e}_x+2\mathbf{e}_y}{2\sqrt{2}}
    = \frac{1}{\sqrt{2}}(\mathbf{e}_x+\mathbf{e}_y).
  \end{equation*}
\end{frame}

\begin{frame}{例：同时求出夹角}
  同一对矢量的点积与模：
  \begin{align*}
    \mathbf{A}\cdot\mathbf{B}
    &= (-1)(1)+1\cdot(-1)+1\cdot 1
     = -1-1+1
     = -1,\\
    |\mathbf{A}|
    &= \sqrt{(-1)^2+1^2+1^2}
     = \sqrt{3},
    \qquad
    |\mathbf{B}|=\sqrt{3}.
  \end{align*}
  所以
  \begin{equation*}
    \cos\theta
    = \frac{-1}{\sqrt{3}\cdot\sqrt{3}}
    = -\frac13,
    \qquad
    \theta
    = \cos^{-1}\!\Bigl(-\frac13\Bigr)
    \approx 109.47^\circ.
  \end{equation*}
  叉积给出法向，点积给出平面内的夹角。
\end{frame}

\subsection{三重积}

\begin{frame}{标量三重积}
  $(\mathbf{A}\times\mathbf{B})\cdot\mathbf{C}$
  与
  $\mathbf{A}\cdot(\mathbf{B}\times\mathbf{C})$
  是同一个标量。
  它的绝对值是三个矢量所张平行六面体的体积，符号记录定向。
  \vspace{0.4em}
  把叉积代入点积：
  \begin{align*}
    \mathbf{A}\cdot(\mathbf{B}\times\mathbf{C})
    &= A_x(B_y C_z-B_z C_y) \\
    &\quad + A_y(B_z C_x-B_x C_z) \\
    &\quad + A_z(B_x C_y-B_y C_x).
  \end{align*}
  这正是行列式按第一行的展开，所以
  \begin{equation*}
    \mathbf{A}\cdot(\mathbf{B}\times\mathbf{C})
    =
    \begin{vmatrix}
      A_x & A_y & A_z \\
      B_x & B_y & B_z \\
      C_x & C_y & C_z
    \end{vmatrix}.
  \end{equation*}
\end{frame}

\begin{frame}{点与叉可以交换}
  把 $(\mathbf{A}\times\mathbf{B})\cdot\mathbf{C}$ 按同样方式展开，
  $x$ 分量贡献
  $(A_y B_z-A_z B_y)C_x$，
  其余两项类似。
  重新集项后得到
  \begin{equation*}
    (\mathbf{A}\times\mathbf{B})\cdot\mathbf{C}
    = A_x(B_y C_z-B_z C_y)+\cdots
    = \mathbf{A}\cdot(\mathbf{B}\times\mathbf{C}).
  \end{equation*}
  交换点与叉不改变结果。
  行列式交换两行变号、循环换行不变号，因此若记
  $[\mathbf{A},\mathbf{B},\mathbf{C}]=\mathbf{A}\cdot(\mathbf{B}\times\mathbf{C})$，则
  \begin{align*}
    [\mathbf{A},\mathbf{B},\mathbf{C}]
    &= [\mathbf{B},\mathbf{C},\mathbf{A}]
     = [\mathbf{C},\mathbf{A},\mathbf{B}] \\
    &= -[\mathbf{A},\mathbf{C},\mathbf{B}]
     = -[\mathbf{B},\mathbf{A},\mathbf{C}]
     = -[\mathbf{C},\mathbf{B},\mathbf{A}].
  \end{align*}
\end{frame}

\begin{frame}{例：体积}
  \begin{columns}[c]
    \begin{column}{0.52\textwidth}
      取
      $\mathbf{A}=\mathbf{e}_x+\mathbf{e}_y$，
      $\mathbf{B}=\mathbf{e}_y+\mathbf{e}_z$，
      $\mathbf{C}=\mathbf{e}_x+\mathbf{e}_z$。
      \begin{align*}
        \mathbf{A}\times\mathbf{B}
        &= \mathbf{e}_x(1\cdot 1-0\cdot 1) \\
        &\quad - \mathbf{e}_y(1\cdot 1-0\cdot 0) \\
        &\quad + \mathbf{e}_z(1\cdot 1-1\cdot 0) \\
        &= \mathbf{e}_x-\mathbf{e}_y+\mathbf{e}_z.
      \end{align*}
    \end{column}
    \begin{column}{0.46\textwidth}
      \centering
      \scalebox{0.72}{\vecfig{parallelepiped}}
    \end{column}
  \end{columns}
\end{frame}

\begin{frame}{例：体积（续）}
  \begin{align*}
    (\mathbf{A}\times\mathbf{B})\cdot\mathbf{C}
    &= (\mathbf{e}_x-\mathbf{e}_y+\mathbf{e}_z)
      \cdot(\mathbf{e}_x+\mathbf{e}_z) \\
    &= 1+0+1
     = 2.
  \end{align*}
  有向体积是 $2$，几何体积是 $|2|=2$。
  用行列式直接验算：
  \begin{equation*}
    \begin{vmatrix}
      1 & 1 & 0 \\
      0 & 1 & 1 \\
      1 & 0 & 1
    \end{vmatrix}
    = 1(1-0)-1(0-1)+0
    = 1+1
    = 2.
  \end{equation*}
\end{frame}

\begin{frame}{拉格朗日恒等式}
  \begin{equation*}
    (\mathbf{A}\times\mathbf{B})\cdot(\mathbf{C}\times\mathbf{D})
    = (\mathbf{A}\cdot\mathbf{C})(\mathbf{B}\cdot\mathbf{D})
    - (\mathbf{A}\cdot\mathbf{D})(\mathbf{B}\cdot\mathbf{C}).
  \end{equation*}
  先把两个叉积写成
  \begin{align*}
    (\mathbf{A}\times\mathbf{B})_i
    &= \varepsilon_{ijk} A_j B_k,\\
    (\mathbf{C}\times\mathbf{D})_i
    &= \varepsilon_{ilm} C_l D_m.
  \end{align*}
  它们的点积是
  \begin{equation*}
    \varepsilon_{ijk}\varepsilon_{ilm}\, A_j B_k C_l D_m.
  \end{equation*}
  缩并公式
  $\varepsilon_{ijk}\varepsilon_{ilm}
  =\delta_{jl}\delta_{km}-\delta_{jm}\delta_{kl}$
  将在第三讲证明。这里先用它把上式化完。
\end{frame}

\begin{frame}{拉格朗日恒等式的化简}
  \begin{align*}
    &\varepsilon_{ijk}\varepsilon_{ilm}\, A_j B_k C_l D_m \\
    &\quad = (\delta_{jl}\delta_{km}-\delta_{jm}\delta_{kl})
      A_j B_k C_l D_m \\
    &\quad = A_j B_k C_j D_k
      - A_j B_k C_k D_j \\
    &\quad = (A_j C_j)(B_k D_k)
      - (A_j D_j)(B_k C_k) \\
    &\quad = (\mathbf{A}\cdot\mathbf{C})(\mathbf{B}\cdot\mathbf{D})
      - (\mathbf{A}\cdot\mathbf{D})(\mathbf{B}\cdot\mathbf{C}).
  \end{align*}
  第一步用了缩并公式，后面两步只是把重复指标收成点积。
\end{frame}

\begin{frame}{矢量三重积}
  $\mathbf{A}\times(\mathbf{B}\times\mathbf{C})$ 垂直于 $\mathbf{A}$，
  并落在 $\mathbf{B}$ 与 $\mathbf{C}$ 所张的平面里。
  叉积没有结合律，两个括法给出不同的矢量：
  \begin{equation*}
    \mathbf{A}\times(\mathbf{B}\times\mathbf{C})
    \quad\text{与}\quad
    (\mathbf{A}\times\mathbf{B})\times\mathbf{C}.
  \end{equation*}
  两条化简公式是
  \begin{align*}
    \mathbf{A}\times(\mathbf{B}\times\mathbf{C})
    &= (\mathbf{A}\cdot\mathbf{C})\mathbf{B}
     - (\mathbf{A}\cdot\mathbf{B})\mathbf{C},\\
    (\mathbf{A}\times\mathbf{B})\times\mathbf{C}
    &= (\mathbf{A}\cdot\mathbf{C})\mathbf{B}
     - (\mathbf{B}\cdot\mathbf{C})\mathbf{A}.
  \end{align*}
  嵌套的叉积因此可以化成原来三个矢量的线性组合。
\end{frame}

\begin{frame}{推导：$x$ 分量}
  先算出
  \begin{align*}
    (\mathbf{B}\times\mathbf{C})_y
    &= B_z C_x-B_x C_z,\\
    (\mathbf{B}\times\mathbf{C})_z
    &= B_x C_y-B_y C_x.
  \end{align*}
  于是
  \begin{align*}
    \bigl[\mathbf{A}\times(\mathbf{B}\times\mathbf{C})\bigr]_x
    &= A_y(B_x C_y-B_y C_x)
     - A_z(B_z C_x-B_x C_z) \\
    &= B_x(A_y C_y+A_z C_z)
     - C_x(A_y B_y+A_z B_z) \\
    &= B_x\bigl(\mathbf{A}\cdot\mathbf{C}-A_x C_x\bigr)
     - C_x\bigl(\mathbf{A}\cdot\mathbf{B}-A_x B_x\bigr) \\
    &= (\mathbf{A}\cdot\mathbf{C})B_x
     - (\mathbf{A}\cdot\mathbf{B})C_x.
  \end{align*}
  $y$、$z$ 分量的算法相同，合起来就是第一条化简公式。
  把公式中的矢量循环置换，得到第二条。
\end{frame}

\begin{frame}{循环和为零}
  把三条化简公式相加：
  \begin{align*}
    &\mathbf{A}\times(\mathbf{B}\times\mathbf{C})
     + \mathbf{B}\times(\mathbf{C}\times\mathbf{A})
     + \mathbf{C}\times(\mathbf{A}\times\mathbf{B}) \\
    &\quad =
      \bigl[(\mathbf{A}\cdot\mathbf{C})\mathbf{B}-(\mathbf{A}\cdot\mathbf{B})\mathbf{C}\bigr] \\
    &\qquad +
      \bigl[(\mathbf{B}\cdot\mathbf{A})\mathbf{C}-(\mathbf{B}\cdot\mathbf{C})\mathbf{A}\bigr] \\
    &\qquad +
      \bigl[(\mathbf{C}\cdot\mathbf{B})\mathbf{A}-(\mathbf{C}\cdot\mathbf{A})\mathbf{B}\bigr].
  \end{align*}
  $\mathbf{A}$、$\mathbf{B}$、$\mathbf{C}$ 的系数分别是
  $-(\mathbf{B}\cdot\mathbf{C})+(\mathbf{C}\cdot\mathbf{B})=0$，
  其余两项同样抵消，所以循环和是 $\mathbf{0}$。
\end{frame}

\begin{frame}{例：垂直于给定方向的分量}
  设 $\hat{\mathbf{n}}$ 是单位矢量，$\hat{\mathbf{n}}\cdot\hat{\mathbf{n}}=1$。
  代入第一条恒等式：
  \begin{align*}
    \hat{\mathbf{n}}\times(\mathbf{v}\times\hat{\mathbf{n}})
    &= (\hat{\mathbf{n}}\cdot\hat{\mathbf{n}})\mathbf{v}
     - (\hat{\mathbf{n}}\cdot\mathbf{v})\hat{\mathbf{n}} \\
    &= \mathbf{v}-(\mathbf{v}\cdot\hat{\mathbf{n}})\hat{\mathbf{n}}.
  \end{align*}
  减去的那一项是 $\mathbf{v}$ 沿 $\hat{\mathbf{n}}$ 的投影，
  剩下的就是垂直于 $\hat{\mathbf{n}}$ 的分量。
\end{frame}

\begin{frame}{例：垂直分量的数值}
  \begin{columns}[c]
    \begin{column}{0.50\textwidth}
      取 $\mathbf{v}=2\mathbf{e}_x-\mathbf{e}_y+3\mathbf{e}_z$，
      $\hat{\mathbf{n}}=\mathbf{e}_z$。
      \begin{align*}
        \mathbf{v}\cdot\mathbf{e}_z &= 3,\\
        \hat{\mathbf{n}}\times(\mathbf{v}\times\hat{\mathbf{n}})
        &= \mathbf{v}-3\mathbf{e}_z \\
        &= 2\mathbf{e}_x-\mathbf{e}_y.
      \end{align*}
      沿选定方向分解横向与纵向分量时就用这个形式。
    \end{column}
    \begin{column}{0.48\textwidth}
      \centering
      \scalebox{0.82}{\vecfig{perp_component}}
    \end{column}
  \end{columns}
\end{frame}

\subsection{从矢量到矩阵}

\begin{frame}{若干矢量排成矩阵}
  把矢量当作列：
  \begin{equation*}
    M
    =
    \begin{pmatrix}
      a_{11} & a_{12} & a_{13} \\
      a_{21} & a_{22} & a_{23} \\
      a_{31} & a_{32} & a_{33}
    \end{pmatrix}.
  \end{equation*}
  矩阵作用在矢量上，就是列的线性组合：
  \begin{equation*}
    M\mathbf{x}
    = x_1\mathbf{a}_1+x_2\mathbf{a}_2+x_3\mathbf{a}_3.
  \end{equation*}
  因此矩阵表示线性变换：它可以旋转、拉伸、反射或剪切。
  用分量写，重复指标 $j$ 要求和，
  \begin{equation*}
    (M\mathbf{x})_i=M_{ij}x_j.
  \end{equation*}
\end{frame}

\begin{frame}{相继两次映射}
  先作用 $N$，再作用 $M$：
  \begin{equation*}
    \bigl(M(N\mathbf{x})\bigr)_i
    = M_{ij}(N\mathbf{x})_j
    = M_{ij} N_{jk} x_k
    = (MN)_{ik} x_k.
  \end{equation*}
  所以
  \begin{equation*}
    M(N\mathbf{x})=(MN)\mathbf{x}.
  \end{equation*}
  单位矩阵的列就是基矢，因而 $I\mathbf{x}=\mathbf{x}$。
\end{frame}

\begin{frame}{单位正方形被送到平行四边形}
  \begin{columns}[c]
    \begin{column}{0.38\textwidth}
      \begin{equation*}
        A=\begin{pmatrix}2&1\\1&2\end{pmatrix}
      \end{equation*}
      的两列是 $\mathbf{e}_x$ 与 $\mathbf{e}_y$ 的像。
      单位正方形被送到这两列所张的平行四边形。
      下一讲把这个映射反过来。
    \end{column}
    \begin{column}{0.60\textwidth}
      \centering
      \scalebox{0.78}{\vecfig{linear_map}}
    \end{column}
  \end{columns}
\end{frame}

\begin{frame}{小结}
  \begin{itemize}
    \item $|\mathbf{A}\times\mathbf{B}|=|\mathbf{A}|\,|\mathbf{B}|\sin\theta$，方向由右手定则给出。
    \item 标量三重积等于三行行列式，是有向体积；轮换不变号，对换变号。
    \item 拉格朗日恒等式把两个叉积的点积化成四个点积。
    \item 矢量三重积化成 $(\mathbf{A}\cdot\mathbf{C})\mathbf{B}-(\mathbf{A}\cdot\mathbf{B})\mathbf{C}$，循环和为零。
    \item 矩阵按列做线性组合；两次映射对应于矩阵相乘。
  \end{itemize}
\end{frame}

\end{document}
